\chapter{运行维护与优化方向}

本章面向现有交付版本，说明运行中需要持续维护的数据、模型配置、工具接口和测试资产。
性能数值和功能核验结果以第七章的测试记录为准，运行维护沿用相同的数据口径和结果结构。

\section{运行维护重点}

\textbf{数据与来源更新。}
任务二写入知识图谱时同时保存数据集标识、文件标识、来源名称和原文片段。
更新公开知识或接入新的 DataMate 数据集时，先登记数据来源和版本，再执行知识构建；
分析库刷新完成后检查图谱事实数量、来源分项和疾病分析表记录是否一致。

\textbf{词典与规则维护。}
医学术语词典、关系规则、否定词规则和文本噪声规则均为独立资产。
新增词条或规则时保留来源、适用实体类型和更新时间，
并使用固定样本检查长词匹配、实体边界、否定表达和结构化字段清洗结果。
离线抽取结果的可复现性只适用于相同输入、规则资产和版本配置，
混合方式的模型结果仍需记录模型配置和调用状态。

\textbf{查询工具与智能体配置。}
Nexent 中的职责提示、约束提示和工具绑定应与 MCP 服务的参数定义同步维护。
工具新增或字段变更后，先运行工具扫描和固定路由用例，
再发布新的智能体版本；任务三由统一分析服务承接查询，继续使用只读连接、允许对象清单、单条语句和返回行数限制。

\section{质量与性能维护规范}

\textbf{抽取质量。}
以 CMeEE、CMeIE 固定划分作为离线比较口径，
分别记录实体和关系的精确率、召回率与 F1，
按同义词缺失、实体边界、类型映射、关系触发词和否定表达分类整理错误样本。
词典和规则每次变更只处理一类问题，并保留变更前后的评测结果，
以便判断改动是否改善了目标类型。

\textbf{知识库更新性能。}
对增量数据采用数据集级任务标识和记录级来源标识，
只处理新增或变更文件；抽取完成后再执行一次性事务写入和分析表刷新。
当数据规模继续增长时，优先检查疾病名称、实体名称、来源标识和关系类型字段的索引，
并用固定数据快照比较单批处理耗时。

\textbf{查询与导出性能。}
继续记录语义规划、SQL 校验、数据库查询、结果组织和文件导出的分阶段耗时。
在线查询采用固定问题集、多轮重复和相同数据库快照，
同时保存中位数、P95（第 95 百分位响应时间）、返回行数和失败次数；
当单次查询接近 5 秒中止限制或返回量接近 200 行上限时，
再针对查询模板、索引和图表截取规则进行调整。
